\section{特性根が非実数の3項間漸化式}\label{節：非実数の特性根}
\subsection{複素数の冪から実数列へ}

第2章では
\[
 a_{n+2}=s a_{n+1}+t a_n
\]
を,\,特性方程式$X^2-sX-t=0$で調べた.
係数と初期二項が実数でも,\,特性根が非実数になる場合がある.
この場合も二根は異なり,$\lambda,\overline\lambda$という共役の組になる.

虚部が正の根を
\[
 \lambda=\rho e^{i\theta},
 \qquad \rho>0,\quad 0<\theta<\pi
\]
と書く.
二つの冪を用いれば
\[
 a_n=A\lambda^{n-1}+B\overline\lambda^{\,n-1}
\]
である.
根が異なるので,\,初めの二項から$A,B$は一意に決まる.
元の各項は実数だから,\,共役を取った式も同じ数列を表す.
係数の一意性から$B=\overline A$となる.

実数$C,D$を使って
\[
 A=\frac{C-iD}{2},\qquad B=\frac{C+iD}{2}
\]
と書き,$u=(n-1)\theta$と置く.
二つの共役な冪の虚部が打ち消し合い,
\[
 \begin{aligned}
 a_n
 &=\rho^{n-1}\{(A+B)\cos u+i(A-B)\sin u\}\\
 &=\rho^{n-1}(C\cos u+D\sin u)
 \end{aligned}
\]
となる.
従って
\[
 \begin{aligned}
 a_n=\rho^{n-1}\{&
 C\cos((n-1)\theta)\\
 &+D\sin((n-1)\theta)\}.
 \end{aligned}
\]
初期条件から
\[
 C=a_1,\qquad
 D=\cfrac{\cfrac{a_2}{\rho}-a_1\cos\theta}{\sin\theta}.
\]
$0<\theta<\pi$なので$\sin\theta\ne0$であり,\,必ず決められる.

三角関数は,\,実数性を保つために組み合わせた二つの複素数の冪を表している.
「複素数を使ったら最後に偶然実数になる」のではなく,\,共役の対称性によって実数になるのである.

\noindent\textbf{【例題】}\par\nobreak
\[
 a_1=0,\quad a_2=1,\qquad
 a_{n+2}=2a_{n+1}-2a_n
\]
の一般項を求めよ.

\noindent\textbf{【方針】}\par\nobreak
特性根$1\pm i$を極形式で表します.

\noindent\textbf{【解答】}\par\nobreak
\[
 a_n=(\sqrt2)^{n-1}\sin\frac{(n-1)\pi}{4}.
\]
\noindent\textbf{【解法】}\par\nobreak
特性方程式は$X^2-2X+2=0$である.
$1+i=\sqrt2e^{i\frac{\pi}{4}}$なので,
$\rho=\sqrt2,\,\theta=\frac{\pi}{4}$となる.
$a_1=0$から$C=0$,
$a_2=1$から$1=\sqrt2D\sin\frac{\pi}{4}=D$と決まる.
最初の数項は$0,1,2,2,0,-4,\ldots$であり,\,得た式とも一致する.
証明は数項の一致ではなく,\,初期二項と更新の一致による.

\subsection{回転を直接記録する複素数}\label{節：複素数の状態}

前の計算では解を冪に分けた.
今度は,\,元の二項から回転を表す複素数を直接作ろう.
根と係数の関係から
\[
 s=\lambda+\overline\lambda,\qquad
 t=-\lambda\overline\lambda.
\]
そこで
\[
 z_n=a_{n+1}-\overline\lambda a_n
\]
と置けば,
\[
 \begin{aligned}
 z_{n+1}
 &=a_{n+2}-\overline\lambda a_{n+1}\\
 &=\lambda a_{n+1}-\lambda\overline\lambda a_n\\
 &=\lambda z_n.
 \end{aligned}
\]
実数列から作った$z_n$は,\,毎回$\rho$倍されながら$\theta$回転する.
第2章で作った補助数列$a_{n+1}-\lambda_1a_n$を,\,複素数の範囲まで広げたのである.

$\lambda=u+iv$で$v>0$と書くと,
\[
 z_n=(a_{n+1}-ua_n)+iv a_n.
\]
従って
\[
 a_n=\frac{\operatorname{Im}z_n}{v},\qquad
 a_{n+1}=\operatorname{Re}z_n+u a_n.
\]
二つの実数と一つの複素数は一対一に対応し,\,情報は失われていない.
元の状態$(a_{n+1},a_n)$から座標を変えたら,\,回転と拡大・縮小が見えたのである.

例えば先ほどの式では$u=v=1$だから
\[
 z_n=(a_{n+1}-a_n)+ia_n.
\]
$z_1=1$より$z_n=(1+i)^{n-1}$となり,\,虚部を取れば一般項が得られる.
\NewFigSpiral

\subsection{元の座標では円とは限らない}

$z_n$の平面では,$\rho=1$なら原点からの距離が一定になる.
しかし元の$(a_{n+1},a_n)$の平面では,\,座標変換が一般に長さや直角を保たないため,\,円上にあるとは限らない.
対応する曲線は
\[
 (a_{n+1}-ua_n)^2+v^2a_n^2=\text{一定}
\]
という楕円になる.
更新が通常の座標でそのまま回転行列である,\,という意味ではない.
実際の状態はその楕円上の点の一部をたどる.

$\rho\ne1$では上の量が毎回$\rho^2$倍される.
この量は,\,元の座標で見た回転の半径に相当する.
固有値が非実数のときに実固有方向が見つからなかったのは,\,固定した一本の実方向にとどまらず,\,平面を回っているからである.

\subsection{周期と収束}

$0<\rho<1$なら,\,正弦と余弦は有界なので$a_n\to0$となる.
$\rho=1$なら有界な振動となるが,\,必ず周期的とは限らない.
非零の実数解が周期$T$を持つための条件は
\[
 \rho=1,\qquad T\theta\in2\pi\mathbb Z
\]
である. ここで$\mathbb Z$は整数全体を表す.

十分性は$\lambda^T=1$から分かる.
必要性は,\,周期があると状態から作った$z_n$にも$z_{n+T}=z_n$が成り立つことから分かる.
非零解では$z_1\ne0$なので$\lambda^Tz_1=z_1$から$\lambda^T=1$となる.
従って周期を持つのは,$\rho=1$かつ$\frac{\theta}{2\pi}$が有理数のときである.
零数列は別に扱う.

$\rho>1$なら非零解の状態は拡大する.
ただし$a_n$だけには零になる項もあり,\,各項の絶対値が単調に増えるとは限らない.
先ほどの例でも4項ごとに零が現れる.
\textbf{状態の大きさと,\,その一成分の大きさを区別する}.
